\documentclass{ximera}

\title{Exercises for Functions and their Representations} 
\license{CC BY-NC-SA 4.0}

\begin{document}

\begin{abstract}
\end{abstract}
\maketitle


%In Exercises \ref{mappingfirst} - \ref{mappinglast}, d



\begin{problem}\label{mappingfirst}
Determine whether or not the mapping diagram represents a function. Explain your reasoning. If the mapping does represent a function, state the domain, range, and represent the function as a set of ordered pairs.

\begin{tikzpicture}[>=Stealth]

% Left labels
\node at (-4.25,5) {Tennant};
\node at (-4.25,3) {Smith};
\node at (-4.25,1) {Calpadi};

% Center label
\node at (0,6) {$M$};

% Right labels
\node at (3.5,5) {Eleven};
\node at (3.5,3) {Twelve};
\node at (3.5,1) {Thirteen};
\node at (3.5,-1) {Fourteen};

% Arrows
\draw[->] (-2.5,5) -- (2.5,5);
\draw[->] (-2.5,5) -- (2.5,3);
\draw[->] (-2.5,3) -- (2.5,1);
\draw[->] (-2.5,1) -- (2.5,-1);

\end{tikzpicture}

\begin{solution}
The mapping $M$ is not a function since  `Tennant' is matched with both `Eleven' and `Twelve.'
\end{solution}
\end{problem} 



\begin{problem}\label{mappinglast}
Determine whether or not the mapping diagram represents a function. Explain your reasoning. If the mapping does represent a function, state the domain, range, and represent the function as a set of ordered pairs.

\begin{tikzpicture}[>=Stealth]

% Left labels
\node at (-4.25,5) {Hartnell};
\node at (-4.25,3) {Cushing};
\node at (-4.25,1) {Hurndall};
\node at (-4.25,-1) {Troughton};

% Center label
\node at (0,6) {$C$};

% Right labels
\node at (3.5,5) {One};
\node at (3.5,3) {Two};
\node at (3.5,1) {Three};

% Arrows
\draw[->] (-2.5,5) -- (2.5,5);
\draw[->] (-2.5,3) -- (2.5,5);
\draw[->] (-2.5,1) -- (2.5,5);
\draw[->] (-2.5,-1) -- (2.5,3);

\end{tikzpicture}

\begin{solution}
The mapping $C$ is a function since each input is matched with only one output.  The domain of $C$ is $\{ \text{Hartnell}, \text{Cushing}, \text{Hurndall}, \text{Troughton} \}$ and the range is $\{\text{One}, \text{Two} \}$.  We can represent $C$ as the following set of ordered pairs:  $\{ (\text{Hartnell}, \text{One}), (\text{Cushing}, \text{One}),  (\text{Hurndall}, \text{One}), (\text{Troughton}, \text{Two}) \}$
\end{solution}
\end{problem} 

\begin{question}
    
In Exercises \ref{tablefirst} - \ref{tablelast}, determine whether or not the data in the given table represents $y$ as a function of $x$.  Explain your reasoning.  If the mapping does represent a function, state the domain, range, and represent the function as a set of ordered pairs.


\begin{problem}

\label{tablefirst}
    \[\begin{array}{|r||r|}  \hline

x  &  y  \\ \hline
 -3 &  3 \\  \hline
 -2 & 2  \\  \hline
  -1 &  1  \\  \hline
 0 &  0 \\  \hline
 1 & 1  \\  \hline
 2 &  2 \\  \hline
 3 & 3  \\  \hline

\end{array}\]

\begin{solution}
In this case, $y$ is a function of $x$ since each $x$ is matched with only one $y$.  

The domain is $\{ -3, -2, -1,0,1,2,3 \}$ and the range is $\{ 0,1,2,3 \}$.  

As ordered pairs, this function is $\{ (-3,3), (-2,2), (-1,1), (0,0), (1,1), (2,2), (3,3) \}$
\end{solution}
\end{problem}

\begin{problem}\label{tablelast}

    \[\begin{array}{|r||r|}  \hline

 x  & y  \\ \hline

 0 & 0 \\  \hline
 1 & 1  \\  \hline
 1 & -1  \\  \hline
 2 &  2 \\  \hline
 2 & -2  \\  \hline
 3 &  3 \\  \hline
 3 & -3  \\  \hline

\end{array}\]

\begin{solution}
In this case, $y$ is not a function of $x$ since there are $x$ values matched with more than one $y$ value.  For instance, $1$ is matched both to $1$ and $-1$.
\end{solution}
\end{problem}

\end{question}




%\setcounter{HW}{\value{enumi}}




%\begin{enumerate}

%\setcounter{enumi}{\value{HW}}

\begin{problem} 
Suppose $W$ is the set of words in the English language and we set up a mapping from $W$ into the set of natural numbers $\mathbb{N}$ as follows: $\text{word} \rightarrow \text{number of letters in the word}$.  Explain why this mapping is a function.  What would you need to know to determine the range of the function?

The mapping is a function since given any word, there is only one answer to `how many letters are in the word?'  For the range, we would need to know what the length of the longest word is and whether or not we could find words of all the lengths between $1$ (the length of the word `a') and it.  See \href{https://en.wikipedia.org/wiki/Longest_word_in_English}{\underline{here}}.
\end{problem}

\begin{problem}
    Suppose $L$ is the set of last names of all the people who have served or are currently serving as the President of the United States.   Consider the mapping from $L$ into $\mathbb{N}$ as follows:  $\text{last name} \rightarrow \text{number of their presidency}$.  For example,  $\text{Washington} \rightarrow 1$ and $\text{Obama} \rightarrow 44$.  Is this mapping a function?  What if we use full names instead of just last names?
    \begin{hint}
    Research what Grover Cleveland and Donald Trump have in common.
    \end{hint}

Since Grover Cleveland was both the 22nd and 24th POTUS, neither mapping described in this exercise is a function.  Donald Trump would also have two ``outputs''.

\end{problem} 

\begin{problem}
  Under what conditions would the time of day be a function of the outdoor temperature? 

\begin{solution}
The outdoor temperature could never be the same for more than two different times - so, for example, it could always be getting warmer or it could always be getting colder.
\end{solution}
\end{problem}

\begin{question}

For the functions $f$ described in Exercises \ref{buildfunctionfirst} - \ref{buildfunctionlast}, find $f(2)$, and find and simplify an expression for $f(x)$ that takes a real number $x$ and performs the following three steps in the order given.


\begin{problem}\label{buildfunctionfirst}
(1) multiply by 2; (2) add 3; (3) divide by 4. 

\begin{solution}
$f(2) = \frac{7}{4}$, $f(x) = \frac{2x+3}{4}$
\end{solution}
\end{problem}


\begin{problem}
(1) add 3; (2) multiply by 2; (3) divide by 4. 

\begin{solution}
$f(2) = \frac{5}{2}$, $f(x) = \frac{2(x+3)}{4} = \frac{x+3}{2}$ 
\end{solution}
\end{problem}

\begin{problem}
(1) divide by 4; (2) add 3; (3) multiply by 2.

\begin{solution}
$f(2) = 7$, $f(x) = 2\left(\frac{x}{4} + 3\right) = \frac{1}{2} x + 6$  
\end{solution}
\end{problem}


\begin{problem}
(1) multiply by 2; (2) add 3; (3) take the square root.

\begin{solution}
$f(2) = \sqrt{7}$, $f(x) = \sqrt{2x+3}$ 
\end{solution}
\end{problem}


\begin{problem}
(1) add 3; (2) multiply by 2; (3) take the square root.

\begin{solution}
$f(2) = \sqrt{10}$, $f(x) = \sqrt{2(x+3)} = \sqrt{2x+6}$
\end{solution}
\end{problem}

\begin{problem}\label{buildfunctionlast}
(1) add 3; (2) take the square root; (3) multiply by 2.

\begin{solution}
$f(2) = 2 \sqrt{5}$, $f(x) = 2\sqrt{x+3}$
\end{solution}
\end{problem}
    
\end{question}


\begin{question}
In Exercises \ref{funcnotationbasicfirst} - \ref{funcnotationbasiclast}, use the given function $f$ to find and simplify the following:

\begin{itemize}
\item $f(3)$
\item $f(-1)$
\item $f\left(\frac{3}{2} \right)$
\item  $f(4x)$
\item $4f(x)$
\item $f(-x)$
\item  $f(x-4)$
\item $f(x) - 4$
\item  $f\left(x^2\right)$
\end{itemize}


\begin{problem}\label{funcnotationbasicfirst} 
$f(x) = 2x+1$

\begin{itemize}
\item $f(3) = \answer{7}$
\item $f(-1) = \answer{-1}$
\item $f\left(\frac{3}{2} \right) = \answer{4}$
\item  $f(4x) = \answer{8x+1}$
\item $4f(x) = \answer{8x+4}$
\item $f(-x) = \answer{-2x+1}$
\item  $f(x-4) = \answer{2x-7}$
\item $f(x) - 4 = \answer{2x-3}$
\item  $f\left(x^2\right) = \answer{2x^2+1}$
\end{itemize}

\end{problem}  

\begin{problem}
$f(x) = 3 - 4x$

\begin{solution}
\item For $f(x) = 3-4x$ 


\begin{itemize}
\item $f(3) = -9$
\item $f(-1) = 7$
\item $f\left(\frac{3}{2} \right) = -3$
\item  $f(4x) = 3-16x$
\item $4f(x) = 12-16x$
\item $f(-x) = 4x+3$
\item  $f(x-4) = 19-4x$
\item $f(x) - 4 = -4x-1$
\item  $f\left(x^2\right) = 3-4x^2$
\end{itemize}

\end{solution}

\end{problem}  

\begin{problem}
$f(x) = 2 - x^2$

\begin{itemize}
\item $f(3) = \answer{-7}$
\item $f(-1) = \answer{1}$
\item $f\left(\frac{3}{2} \right) = \answer{-\frac{1}{4}}$
\item  $f(4x) = \answer{2-16x^2}$
\item $4f(x) = \answer{8-4x^2}$
\item $f(-x) = \answer{2-x^2}$
\item  $f(x-4) = \answer{-x^2+8x-14}$
\item $f(x) - 4 = \answer{-x^2-2}$
\item  $f\left(x^2\right) = \answer{2-x^4}$
\end{itemize}

\end{problem}  


\begin{problem}
$f(x) = x^2 - 3x + 2$,


\begin{itemize}
\item $f(3) = \answer{2}$
\item $f(-1) = \answer{6}$
\item $f\left(\frac{3}{2} \right) = \answer{-\frac{1}{4}}$
\item  $f(4x) = \answer{16x^2-12x+2}$
\item $4f(x) = \answer{4x^2-12x+8}$
\item $f(-x) = \answer{x^2+3x+2}$
\item  $f(x-4) = \answer{x^2-11x+30}$
\item $f(x) - 4 = \answer{x^2-3x-2}$
\item  $f\left(x^2\right) = \answer{x^4-3x^2+2}$
\end{itemize}  
\end{problem}

\begin{problem}
$f(x) = 6$

\begin{itemize}
\item $f(3) = \answer{6}$
\item $f(-1) =\answer{6}$
\item $f\left(\frac{3}{2} \right) = \answer{6}$
\item  $f(4x) = \answer{6}$
\item $4f(x) = \answer{24}$
\item $f(-x) = \answer{6}$
\item  $f(x-4) = \answer{6}$ 
\item $f(x) - 4 = \answer{2}$
\item  $f\left(x^2\right) = \answer{6}$
\end{itemize}

\end{problem}  

\begin{problem}\label{funcnotationbasiclast}
$f(x) = 0$

\begin{itemize}
\item $f(3) = \answer{0}$
\item $f(-1) =\answer{0}$
\item $f\left(\frac{3}{2} \right) = \answer{0}$
\item  $f(4x) = \answer{0}$
\item $4f(x) = \answer{0}$
\item $f(-x) = \answer{0}$
\item  $f(x-4) = \answer{0}$ 
\item $f(x) - 4 = \answer{-4}$
\item  $f\left(x^2\right) = \answer{0}$
\end{itemize}

\end{problem}  

\end{question}

\begin{question}
In Exercises \ref{secondfuncnotationbasicfirst} - \ref{secondfuncnotationbasiclast}, use the given function $f$ to find and simplify the following:

\begin{itemize}
\item  $f(2)$
\item  $f(-2)$
\item  $f(2a)$
\item  $2 f(a)$
\item $f(a+2)$
\item $f(a) + f(2)$
\item  $f \left( \frac{2}{a} \right)$
\item $\frac{f(a)}{2}$
\item  $f(a + h)$
\end{itemize}
    
\begin{problem}\label{secondfuncnotationbasicfirst}
$f(x) = 2x-5$

\begin{itemize}
\item  $f(2)=\answer{-1}$
\item  $f(-2)=\answer{-9}$
\item  $f(2a)=\answer{4a-5}$
\item  $2 f(a)=\answer{4a-10}$
\item $f(a+2)=\answer{2a-1}$
\item $f(a) + f(2)=\answer{2a-6}$
\item  $f \left( \frac{2}{a} \right)=\answer{\frac{4}{a} - 5}$
\item $\frac{f(a)}{2}=\answer{\frac{2a-5}{2}}$
\item  $f(a + h)=\answer{2a + 2h - 5}$
\end{itemize}

\end{problem}

\begin{problem}
$f(t) = 5-2t$

\begin{itemize}
\item  $f(2)=\answer{1}$
\item  $f(-2)=\answer{9}$
\item  $f(2a)=\answer{5-4a}$
\item  $2 f(a)=\answer{10-4a}$
\item $f(a+2)=\answer{1-2a}$
\item $f(a) + f(2)=\answer{6-2a}$
\item  $f \left( \frac{2}{a} \right)=\answer{5 - \frac{4}{a}}$
\item $\frac{f(a)}{2}=\answer{\frac{5-2a}{2}}$
\item  $f(a + h)=\answer{5-2a-2h}$
\end{itemize}

\end{problem}


\begin{problem}
$f(w) = 2w^2 - 1$

\begin{itemize}
\item  $f(2)=\answer{7}$
\item  $f(-2)=\answer{7}$
\item  $f(2a)=\answer{8a^2-1}$
\item  $2 f(a)=\answer{4a^2-2}$
\item $f(a+2)=\answer{2a^2+8a+7}$
\item $f(a) + f(2)=\answer{2a^2+6}$
\item  $f \left( \frac{2}{a} \right)=\answer{\frac{8}{a^2} - 1}$
\item $\frac{f(a)}{2}=\answer{\frac{2a^2-1}{2}}$
\item  $f(a + h)=\answer{2a^2+4ah+2h^2-1}$
\end{itemize}

\end{problem}


\begin{problem}
$f(q) = 3q^2+3q-2$


\begin{itemize}
\item  $f(2)=\answer{16}$
\item  $f(-2)=\answer{4}$
\item  $f(2a)=\answer{12a^2+6a-2}$
\item  $2 f(a)=\answer{6a^2+6a-4}$
\item $f(a+2)=\answer{3a^2+15a+16}$
\item $f(a) + f(2)=\answer{3a^2+3a+14}$
\item  $f \left( \frac{2}{a} \right)=\answer{\frac{12}{a^2} + \frac{6}{a} - 2}$
\item $\frac{f(a)}{2}=\answer{\frac{3a^2+3a-2}{2}}$
\item  $f(a + h)=\answer{3a^2 + 6ah + 3h^2+3a+3h-2}$
\end{itemize}

\end{problem}


\begin{problem}
$f(r) = 117$

\begin{itemize}
\item  $f(2)=\answer{117}$
\item  $f(-2)=\answer{117}$
\item  $f(2a)=\answer{117}$
\item  $2 f(a)=\answer{234}$
\item $f(a+2)=\answer{117}$
\item $f(a) + f(2)=\answer{117}$
\item  $f \left( \frac{2}{a} \right)=\answer{117}$
\item $\frac{f(a)}{2}=\answer{\frac{117}{2}}$
\item  $f(a + h)=\answer{117}$
\end{itemize}

\end{problem}
  

\begin{problem}\label{secondfuncnotationbasiclast}
$f(z) = \frac{z}{2}$

\begin{itemize}
\item  $f(2)=\answer{1}$
\item  $f(-2)=\answer{-1}$
\item  $f(2a)=\answer{a}$
\item  $2 f(a)=\answer{a}$
\item $f(a+2)=\answer{\frac{a+2}{2}}$
\item $f(a) + f(2)=\answer{\frac{a}{2}+ 1}$
\item  $f \left( \frac{2}{a} \right)=\answer{\frac{1}{a}}$
\item $\frac{f(a)}{2}=\answer{\frac{a}{4}}$
\item  $f(a + h)=\answer{\frac{a+h}{2}}$
\end{itemize}

\end{problem}
\end{question}

\begin{question}
In Exercises \ref{findzerofuncfirst} - \ref{findzerofunclast}, use the given function $f$ to find $f(0)$ and solve $f(x)=0$.

\begin{problem}\label{findzerofuncfirst}
$f(x) = 2x - 1$

$f(0) = \answer{-1}$

Solving $f(x)=0$ gives us 
\begin{selectAll}
    \choice{$x=-1$}
    \choice{$x=-\frac{1}{2}$}
    \choice{$x=0$}
    \choice[correct]{$x=\frac{1}{2}$}
    \choice{$x=1$}
\end{selectAll}
\end{problem}


\begin{problem}
$f(x) = 3 - \frac{2}{5} x$

$f(0) = \answer{3}$

Solving $f(x)=0$ gives us 
\begin{selectAll}
    \choice{$x=0$}
    \choice{$x=\frac{6}{5}$}
    \choice{$x=3$}
    \choice{$x=\frac{5}{13}$}
    \choice{$x=5$}
    \choice[correct]{$x=\frac{15}{2}$}
\end{selectAll}
\end{problem}


\begin{problem}
$f(x) = 2x^2 - 6$,

$f(0) = \answer{-6}$

Solving $f(x)=0$ gives us 
\begin{selectAll}
    \choice{$x=-6$}
    \choice{$x=-3$}
    \choice[correct]{$x=-\sqrt{3}$}
    \choice{$x=0$}
    \choice[correct]{$x=\sqrt{3}$}
    \choice{$x=3$}
    \choice{$x=6$}
\end{selectAll}
\end{problem}

\begin{problem}\label{findzerofunclast}
$f(x) = x^2 - x - 12$

$f(0) = \answer{-12}$

Solving $f(x)=0$ gives us 
\begin{selectAll}
    \choice{$x=-12$}
    \choice{$x=-4$}
    \choice[correct]{$x=-3$}
    \choice{$x=0$}
    \choice{$x=3$}
    \choice[correct]{$x=4$}
    \choice{$x=12$}
\end{selectAll}
\end{problem}

\end{question}

\begin{question}
In Exercises \ref{equfunctionfirst} - \ref{equfunctionlast}, determine whether or not the following equation represents $y$ as a function of $x$.

\begin{problem}\label{equfunctionfirst}
$y = x^{3} - x$

\begin{multipleChoice}
    \choice[correct]{Yes, $y$ is a function of $x$.}
    \choice{No, $y$ is a not a function of $x$.}
\end{multipleChoice}

\end{problem}

\begin{problem}
$y = \sqrt{x - 2}$

\begin{multipleChoice}
    \choice[correct]{Yes, $y$ is a function of $x$.}
    \choice{No, $y$ is a not a function of $x$.}
\end{multipleChoice}
\end{problem}

\begin{problem}
$x^{3}y = -4$ 

\begin{multipleChoice}
    \choice[correct]{Yes, $y$ is a function of $x$.}
    \begin{feedback}
        Yes, in fact, $y=\frac{-4}{x^3}$.
    \end{feedback}
    \choice{No, $y$ is a not a function of $x$.}
\end{multipleChoice}
\end{problem}

\begin{problem}
$x^{2} - y^{2} = 1$ 

\begin{multipleChoice}
    \choice{Yes, $y$ is a function of $x$.}
    \choice[correct]{No, $y$ is a not a function of $x$.}
    \begin{feedback}
        We cannot solve for y without using $\pm$ when we take the square root of both sides of our equation.
    \end{feedback}
\end{multipleChoice}
\end{problem}


\begin{problem}
$y = \frac{x}{x^{2} - 9}$ 

\begin{multipleChoice}
    \choice[correct]{Yes, $y$ is a function of $x$.}
    \choice{No, $y$ is a not a function of $x$.}
\end{multipleChoice}
\end{problem}


\begin{problem}
$x = -6$
\begin{multipleChoice}
    \choice{Yes, $y$ is a function of $x$.}
    \choice[correct]{No, $y$ is a not a function of $x$.}
\end{multipleChoice}
\end{problem}


\begin{problem}
$x = y^2 + 4$
\begin{multipleChoice}
    \choice{Yes, $y$ is a function of $x$.}
    \choice[correct]{No, $y$ is a not a function of $x$.}
    \begin{feedback}
        We cannot solve for y without using $\pm$ when we take the square root of both sides of our equation.
    \end{feedback}
\end{multipleChoice}
\end{problem}  

\begin{problem}
$y = x^2 + 4$
\begin{multipleChoice}
    \choice[correct]{Yes, $y$ is a function of $x$.}
    \choice{No, $y$ is a not a function of $x$.}
\end{multipleChoice}
\end{problem}

\begin{problem}
$x^2 + y^2 = 4$
\begin{multipleChoice}
    \choice{Yes, $y$ is a function of $x$.}
    \choice[correct]{No, $y$ is a not a function of $x$.}
    \begin{feedback}
        We cannot solve for y without using $\pm$ when we take the square root of both sides of our equation.
    \end{feedback}
\end{multipleChoice}
\end{problem}

\begin{problem}
$y = \sqrt{4-x^2}$
\begin{multipleChoice}
    \choice[correct]{Yes, $y$ is a function of $x$.}
    \choice{No, $y$ is a not a function of $x$.}
\end{multipleChoice}
\end{problem}

\begin{problem}
$x^2 - y^2 = 4$
\begin{multipleChoice}
    \choice{Yes, $y$ is a function of $x$.}
    \choice[correct]{No, $y$ is a not a function of $x$.}
    \begin{feedback}
        We cannot solve for y without using $\pm$ when we take the square root of both sides of our equation.
    \end{feedback}
\end{multipleChoice}
\end{problem}

\begin{problem}
$x^3 + y^3 = 4$
\begin{multipleChoice}
    \choice[correct]{Yes, $y$ is a function of $x$.}
    \begin{feedback}
        Yes, in fact, $y=\sqrt[3]{4-x^3}$.
    \end{feedback}
    \choice{No, $y$ is a not a function of $x$.}
\end{multipleChoice}
\end{problem}


\begin{problem}
$2x + 3y = 4$
\begin{multipleChoice}
    \choice[correct]{Yes, $y$ is a function of $x$.}
    \begin{feedback}
        Yes, in fact, $y=-\frac{2}{3}x+\frac{4}{3}$.
    \end{feedback}
    \choice{No, $y$ is a not a function of $x$.}
\end{multipleChoice}
\end{problem}

\begin{problem}
$2xy = 4$
\begin{multipleChoice}
    \choice[correct]{Yes, $y$ is a function of $x$.}
    \begin{feedback}
        Yes, in fact, $y=\frac{4}{2x}$.
    \end{feedback}
    \choice{No, $y$ is a not a function of $x$.}
\end{multipleChoice}
\end{problem}

\begin{problem}\label{equfunctionlast}
$x^2 = y^2$
\begin{multipleChoice}
    \choice{Yes, $y$ is a function of $x$.}
    \choice[correct]{No, $y$ is a not a function of $x$.}
    \begin{feedback}
        We cannot solve for y without using $\pm$ when we take the square root of both sides of our equation.
    \end{feedback}
\end{multipleChoice}
\end{problem}

\end{question}

\begin{question}
Exercises \ref{setfunctionfirst} - \ref{setfunctionlast} give a set of points in the $xy$-plane. Determine if $y$ is a function of $x$. If so, state the
domain and range.

\begin{problem}\label{setfunctionfirst}
\{$(-3, 9)$, $\;(-2, 4)$, $\;(-1, 1)$, $\;(0, 0)$, $\;(1, 1)$, $\;(2, 4)$, $\;(3, 9)\}$ 

\begin{multipleChoice}
    \choice[correct]{Yes, $y$ is a function of $x$.}
    \choice{No, $y$ is a not a function of $x$.}
    \begin{feedback}
    \begin{solution}
       domain = \{$-3$, $-2$, $-1$, $0$, $1$, $2$ ,$3$\} \\ range = \{$0$, $1$, $4$, $9$\} 
    \end{solution}
    \end{feedback}
\end{multipleChoice}
\end{problem}

\begin{problem}
$\left\{ (-3,0), (1,6), (2,-3), (4,2), (-5,6), (4,-9), (6,2) \right\}$

\begin{multipleChoice}
    \choice{Yes, $y$ is a function of $x$.}
    \choice[correct]{No, $y$ is a not a function of $x$.}
    \begin{feedback}
       Notice $4$ has two ``outputs''. 
    \end{feedback}
\end{multipleChoice}
\end{problem}   


\begin{problem}
$\left\{ (-3,0), (-7,6), (5,5), (6,4), (4,9), (3,0) \right\}$

\begin{multipleChoice}
    \choice[correct]{Yes, $y$ is a function of $x$.}
    \choice{No, $y$ is a not a function of $x$.}
    \begin{feedback}
    \begin{solution}
       domain = $\left\{ -7, -3, 3, 4, 5, 6 \right\}$ \\ range = $\left\{ 0,4,5,6,9 \right\}$ 
    \end{solution}
    \end{feedback}
\end{multipleChoice}
\end{problem}    

\begin{problem}
$\left\{ (1,2), (4,4), (9,6), (16,8), (25,10), (36, 12), \ldots \right\}$

\begin{multipleChoice}
    \choice[correct]{Yes, $y$ is a function of $x$.}
    \choice{No, $y$ is a not a function of $x$.}
    \begin{feedback}
    \begin{solution}
       domain =   $\left\{ 1, 4, 9, 16, 25, 36, \ldots \right\} \\ = \left\{ x \, | \, \text{$x$ is a perfect square} \right\}$ \\ range =  $\left\{ 2, 4, 6, 8, 10, 12, \ldots \right\} \\ = \left\{ y \, | \, \text{$y$ is a positive even integer} \right\}$ 
    \end{solution}
    \end{feedback}
\end{multipleChoice}
\end{problem}    

\begin{problem}
\{($x, y) \, | \, x$ is an odd integer, and $y$ is an even integer\}

\begin{multipleChoice}
    \choice{Yes, $y$ is a function of $x$.}
    \choice[correct]{No, $y$ is a not a function of $x$.}
    \begin{feedback}
       Notice that each ``input'' has many ``outputs''. 
    \end{feedback}
\end{multipleChoice}
\end{problem}  

\begin{problem}
\{$(x, 1) \, | \, x$ is an irrational number\}

\begin{multipleChoice}
    \choice[correct]{Yes, $y$ is a function of $x$.}
    \choice{No, $y$ is a not a function of $x$.}
    \begin{feedback}
    \begin{solution}
       domain $= \{x \, | \, \text{$x$ is irrational} \}$ \\ range $= \{ 1\}$ 
    \end{solution}
    \end{feedback}
\end{multipleChoice}
\end{problem}   

\begin{problem}
$\{ (1,0), (2,1), (4,2), (8,3), (16,4), (32, 5), \ldots \}$

\begin{multipleChoice}
    \choice[correct]{Yes, $y$ is a function of $x$.}
    \choice{No, $y$ is a not a function of $x$.}
    \begin{feedback}
    \begin{solution}
       domain  $= \{x \, | \, 1, 2, 4, 8, \ldots \}$ \\ $= \{x \, | \, \text{$x = 2^{n}$ for some whole number $n$} \}$ \\ range $= \{ 0, 1, 2, 3, \ldots \}$ \\ $= \{y \, | \, \text{$y$ is any whole number}\}$ 
    \end{solution}
    \end{feedback}
\end{multipleChoice}
\end{problem}   

\begin{problem}
$\{ \ldots (-3,9), (-2,4), (-1,1), (0,0), (1,1), (2,4), (3,9), \ldots \}$

\begin{multipleChoice}
    \choice[correct]{Yes, $y$ is a function of $x$.}
    \choice{No, $y$ is a not a function of $x$.}
    \begin{feedback}
    \begin{solution}
domain $= \{x \, | \, \text{$x$ is any integer} \}$ \\ range $= \{y \, | \, \text{ $y$ is the square of an integer}\}$
    \end{solution}
    \end{feedback}
\end{multipleChoice}
\end{problem}   

\begin{problem}
$\{ (-2, y) \, | \, -3 < y < 4\}$

\begin{multipleChoice}
    \choice{Yes, $y$ is a function of $x$.}
    \choice[correct]{No, $y$ is a not a function of $x$.}
    \begin{feedback}
       Notice $-2$ has many``outputs''. 
    \end{feedback}
\end{multipleChoice}
\end{problem}   

\begin{problem}
$\{ (x,3) \, | \,  -2 \leq x < 4\}$

\begin{multipleChoice}
    \choice[correct]{Yes, $y$ is a function of $x$.}
    \choice{No, $y$ is a not a function of $x$.}
    \begin{feedback}
    \begin{solution}
       domain  $= \{x \, | \, -2 \leq x < 4 \} = [-2, 4)$, \\ range = \{$3$\} 
    \end{solution}
    \end{feedback}
\end{multipleChoice}
\end{problem}    


\begin{problem}
$\{ \left(x,x^2\right) \, | \, \text{$x$ is a real number} \}$

\begin{multipleChoice}
    \choice[correct]{Yes, $y$ is a function of $x$.}
    \choice{No, $y$ is a not a function of $x$.}
    \begin{feedback}
    \begin{solution}
       domain $= \{x \, | \,  \text{$x$ is a real number} \} = (-\infty, \infty)$ \\  range $= \{y \, | \,  y \geq 0 \} = [0,\infty)$ 
    \end{solution}
    \end{feedback}
\end{multipleChoice}
\end{problem}    

\begin{problem}\label{setfunctionlast}
$\{ \left(x^2,x\right) \, | \, \text{$x$ is a real number} \}$

\begin{multipleChoice}
    \choice{Yes, $y$ is a function of $x$.}
    \choice[correct]{No, $y$ is a not a function of $x$.}
    \begin{feedback}
       Notice $(4,2)$ and $(4,-2)$ are both in this set, for example. 
    \end{feedback}
\end{multipleChoice}
\end{problem}    
    
\end{question}
\begin{problem}\label{HLTExercise} The Vertical Line Test is a quick way to determine from a graph if the vertical axis variable is a function of the horizontal axis variable. If we are given a graph and asked to determine if the horizontal axis variable is a function of the vertical axis variable, we can use horizontal lines instead of vertical lines to check.  Using Theorem \ref{VLT} as a guide,  formulate a `Horizontal Line Test.'  (We'll refer back to this exercise in Section \ref{InverseFunctions}.)

\begin{solution}
\textbf{Horizontal Line Test:} A graph on the $xy$-plane represents $x$ as a function of $y$ if and only if no \textbf{horizontal} line intersects the graph more than once.
\end{solution}
\end{problem}

%\setcounter{HW}{\value{enumi}}


 


In Exercises \ref{graphfunctionfirstxy} - \ref{graphfunctionlastxy}, determine whether or not the graph suggests $y$ is a function of $x$.  For the ones which do, state the domain and range. 


\begin{problem}\label{graphfunctionfirstxy}

    \begin{center}
\begin{tikzpicture}
\begin{axis}[
    width=2.7in, height=2.4in,
    xmin=-5, xmax=2, ymin=-2, ymax=5,
    axis lines=middle,
    axis line style={-{Latex[length=8pt,width=6pt]}},
    tick style={line width=1pt},
    major tick length=5pt,
    xlabel={$x$}, ylabel={$y$},
    xlabel style={anchor=west},
    ylabel style={anchor=south},
    xtick={-4,-3,-2,-1,1},
    ytick={-1,1,2,3,4},
    tick label style={font=\tiny},
    label style={font=\scriptsize},
    samples=100,
]
\addplot[only marks, mark=*, mark size=2pt] coordinates {(-4,-1) (-3,0) (-2,1) (-1,2) (0,3) (1,4)};
\end{axis}
\end{tikzpicture}
\end{center}

\begin{multipleChoice}
    \choice[correct]{Yes, $y$ is a function of $x$.}
    \choice{No, $y$ is a not a function of $x$.}
    \begin{feedback}
    \begin{solution}
       domain = \{$-4$, $-3$, $-2$, $-1$, $0$, $1$\} \\ range = \{$-1$, $0$, $1$, $2$, $3$, $4$\} 
    \end{solution}
    \end{feedback}
\end{multipleChoice}
\end{problem}    



%  
%  

\begin{problem}\label{graphfunctionfirstxy2}


    \begin{center}
\begin{tikzpicture}
\begin{axis}[
    width=2.7in, height=2.4in,
    xmin=-5, xmax=2, ymin=-2, ymax=5,
    axis lines=middle,
    axis line style={-{Latex[length=8pt,width=6pt]}},
    tick style={line width=1pt},
    major tick length=5pt,
    xlabel={$x$}, ylabel={$y$},
    xlabel style={anchor=west},
    ylabel style={anchor=south},
    xtick={-4,-3,-2,-1,1},
    ytick={-1,1,2,3,4},
    tick label style={font=\tiny},
    label style={font=\scriptsize},
    samples=100,
]
\addplot[only marks, mark=*, mark size=2pt] coordinates {(-4,-1) (-3,0) (-3,1) (-2,1) (-1,2) (0,3) (1,4)};
\end{axis}
\end{tikzpicture}
\end{center}


\begin{multipleChoice}
    \choice{Yes, $y$ is a function of $x$.}
    \choice[correct]{No, $y$ is a not a function of $x$.}
    \begin{feedback}
       Notice the vertical line test fails for $x=-3$. 
    \end{feedback}
\end{multipleChoice}
\end{problem}  



%\setcounter{HW}{\value{enumi}}

%\setcounter{enumi}{\value{HW}}

\begin{problem}\label{graphfunctionfirstxy3}


    \begin{center}
\begin{tikzpicture}
\begin{axis}[
    width=2.7in, height=2.4in,
    xmin=-3, xmax=3, ymin=-1, ymax=6,
    axis lines=middle,
    axis line style={-{Latex[length=8pt,width=6pt]}},
    tick style={line width=1pt},
    major tick length=5pt,
    xlabel={$x$}, ylabel={$y$},
    xlabel style={anchor=west},
    ylabel style={anchor=south},
    xtick={-2,-1,1,2},
    ytick={1,2,3,4,5},
    tick label style={font=\tiny},
    label style={font=\scriptsize},
    samples=100,
]
\addplot[<->, line width=1.25pt, domain=-2.1:2.1] {x^2+1};
\end{axis}
\end{tikzpicture}
\end{center}
\begin{multipleChoice}
    \choice[correct]{Yes, $y$ is a function of $x$.}
    \choice{No, $y$ is a not a function of $x$.}
    \begin{feedback}
    \begin{solution}
domain = $(-\infty, \infty)$ \\ range = $[1, \infty)$
    \end{solution}
    \end{feedback}
\end{multipleChoice}
\end{problem}


 
\begin{problem}\label{graphfunctionlastxy}

    \begin{center}
\begin{tikzpicture}
\begin{axis}[
    width=2.7in, height=2.4in,
    xmin=-4, xmax=4, ymin=-4, ymax=4,
    axis lines=middle,
    axis line style={-{Latex[length=8pt,width=6pt]}},
    tick style={line width=1pt},
    major tick length=5pt,
    xlabel={$x$}, ylabel={$y$},
    xlabel style={anchor=west},
    ylabel style={anchor=south},
    xtick={-3,-2,-1,1,2,3},
    ytick={-3,-2,-1,1,2,3},
    tick label style={font=\tiny},
    label style={font=\scriptsize},
    samples=100,
]
\addplot[<->, line width=1.25pt, domain=-2:2, samples=100] ({cosh(x)}, {sinh(x)});
\end{axis}
\end{tikzpicture}
\end{center}

\begin{multipleChoice}
    \choice{Yes, $y$ is a function of $x$.}
    \choice[correct]{No, $y$ is a not a function of $x$.}
    \begin{feedback}
       Notice the vertical line test fails many places. 
    \end{feedback}
\end{multipleChoice}
\end{problem} 



%\setcounter{HW}{\value{enumi}}

%\setcounter{enumi}{\value{HW}}

\begin{problem}
    Determine which of the graphs, if any, in Exercises \ref{graphfunctionfirstxy} - \ref{graphfunctionlastxy}, suggest $x$ is a function of $y$.  (Feel free to use Exercise \ref{HLTExercise}.)  If it does, state the domain and range.

\begin{solution}
\begin{itemize}

\item Number \ref{graphfunctionfirstxy} represents $x$ as a function of $y$.

domain  = \{$-1$, $0$, $1$, $2$, $3$, $4$\} and range = \{$-4$, $-3$, $-2$, $-1$, $0$, $1$\}

\item Number \ref{graphfunctionlastxy} represents $x$ as a function of $y$.

domain = $(-\infty, \infty)$  and range = $[1, \infty)$


\end{itemize}
\end{solution}
\end{problem}

  

%\setcounter{HW}{\value{enumi}}



%\setcounter{enumi}{\value{HW}}

\begin{question}
  In Exercises \ref{graphfunctionfirstvw} - \ref{graphfunctionlastvw}, determine whether or not the graph suggests $w$ is a function of $v$.   For the ones which do, state the domain and range.

\begin{problem}\label{graphfunctionfirstvw}


    \begin{center}
\begin{tikzpicture}
\begin{axis}[
    width=2.7in, height=2.4in,
    xmin=-1, xmax=10, ymin=-1, ymax=4,
    axis lines=middle,
    axis line style={-{Latex[length=8pt,width=6pt]}},
    tick style={line width=1pt},
    major tick length=5pt,
    xlabel={$v$}, ylabel={$w$},
    xlabel style={anchor=west},
    ylabel style={anchor=south},
    xtick={1,2,3,4,5,6,7,8,9},
    ytick={1,2,3},
    tick label style={font=\tiny},
    label style={font=\scriptsize},
    samples=100,
]
\addplot[->, line width=1.25pt, domain=2:10] {sqrt(x-2)};
\addplot[only marks, mark=*, mark size=2pt] coordinates {(2,0)};
\end{axis}
\end{tikzpicture}
\end{center}

\begin{multipleChoice}
    \choice[correct]{Yes, $w$ is a function of $v$.}
    \choice{No, $w$ is a not a function of $v$.}
    \begin{feedback}
    \begin{solution}
domain = $[2, \infty)$ \\ range = $[0, \infty)$
    \end{solution}
    \end{feedback}
\end{multipleChoice}
\end{problem}

\begin{problem}\label{graphfunctionfirstvw2}


    \begin{center}
\begin{tikzpicture}
\begin{axis}[
    width=2.7in, height=2.4in,
    xmin=-5, xmax=5, ymin=-1, ymax=5,
    axis lines=middle,
    axis line style={-{Latex[length=8pt,width=6pt]}},
    tick style={line width=1pt},
    major tick length=5pt,
    xlabel={$v$}, ylabel={$w$},
    xlabel style={anchor=west},
    ylabel style={anchor=south},
    xtick={-4,-3,-2,-1,1,2,3,4},
    ytick={1,2,3,4},
    tick label style={font=\tiny},
    label style={font=\scriptsize},
    samples=100,
]
\addplot[<->, line width=1.25pt, domain=-5:5] {4/(x^2+1)};
\end{axis}
\end{tikzpicture}
\end{center}

\begin{multipleChoice}
    \choice[correct]{Yes, $w$ is a function of $v$.}
    \choice{No, $w$ is a not a function of $v$.}
    \begin{feedback}
    \begin{solution}
domain = $[-\infty, \infty)$ \\ range = $(0, \infty)$
    \end{solution}
    \end{feedback}
\end{multipleChoice}
\end{problem}


\begin{problem}\label{graphfunctionfirstvw3}

    \begin{center}
\begin{tikzpicture}
\begin{axis}[
    width=2.7in, height=2.4in,
    xmin=-4.5, xmax=5.5, ymin=-4, ymax=3,
    axis lines=middle,
    axis line style={-{Latex[length=8pt,width=6pt]}},
    tick style={line width=1pt},
    major tick length=5pt,
    xlabel={$v$}, ylabel={$w$},
    xlabel style={anchor=west},
    ylabel style={anchor=south},
    xtick={-4,-2,-1,1,2,4,5},
    ytick={-2,-1,1},
    tick label style={font=\tiny},
    label style={font=\scriptsize},
    samples=100,
]
\addplot[line width=1.25pt] coordinates {(-3,-3) (-3,2) (3,2) (3,-3) (-3,-3)};
\end{axis}
\end{tikzpicture}
\end{center}

\begin{multipleChoice}
    \choice{Yes, $w$ is a function of $v$.}
    \choice[correct]{No, $w$ is a not a function of $v$.}
    \begin{feedback}
    The vertical line test fails in many places.
    \end{feedback}
\end{multipleChoice}
\end{problem}   




\begin{problem}\label{graphfunctionlastvw}

    \begin{center}
\begin{tikzpicture}
\begin{axis}[
    width=2.7in, height=2.4in,
    xmin=-6, xmax=4, ymin=-2.5, ymax=4.5,
    axis lines=middle,
    axis line style={-{Latex[length=8pt,width=6pt]}},
    tick style={line width=1pt},
    major tick length=5pt,
    xlabel={$v$}, ylabel={$w$},
    xlabel style={anchor=west},
    ylabel style={anchor=south},
    xtick={-5,-4,-3,-2,-1,1,2,3},
    ytick={-2,-1,1,2,3,4},
    tick label style={font=\tiny},
    label style={font=\scriptsize},
    samples=100,
]
\addplot[line width=1.25pt, domain=-5:-1] {-5-6*x-x^2};
\addplot[line width=1.25pt, domain=-1:3] {x/4-7/4};
\addplot[only marks, mark=*, mark size=2pt] coordinates {(-5,0) (-1,0)};
\addplot[only marks, mark=*, mark size=2pt, mark options={fill=white}] coordinates {(-3,4) (-1,-2) (3,-1)};
\end{axis}
\end{tikzpicture}
\end{center}
\begin{multipleChoice}
    \choice[correct]{Yes, $w$ is a function of $v$.}
    \choice{No, $w$ is a not a function of $v$.}
    \begin{feedback}
    \begin{solution}
domain = $[-5,-3) \cup(-3, 3)$ \\ range = $(-2, -1) \cup [0, 4)$
    \end{solution}
    \end{feedback}
\end{multipleChoice}
\end{problem}

\end{question}


\begin{problem}
    Determine which of the graphs, if any, in Exercises \ref{graphfunctionfirstvw} - \ref{graphfunctionlastvw}, suggest $v$ is a function of $w$.  (Feel free to use Exercise \ref{HLTExercise}.)  If it does, state the domain and range.

\begin{solution}
Only number \ref{graphfunctionfirstvw} represents $v$ as a function of $w$;  domain = $[0, \infty)$ and range = $[2, \infty)$
\end{solution}
\end{problem}


\begin{question}
In Exercises \ref{graphfunctionfirsttT} - \ref{graphfunctionlasttT}, determine whether or not the graph suggests $T$ is a function of $t$.   For the ones which do, state the domain and range. 



%\setcounter{enumi}{\value{HW}}

\begin{problem}\label{graphfunctionfirsttT}


    \begin{center}
\begin{tikzpicture}
\begin{axis}[
    width=2.7in, height=3in,
    xmin=-4, xmax=4, ymin=-6, ymax=10,
    axis lines=middle,
    axis line style={-{Latex[length=8pt,width=6pt]}},
    tick style={line width=1pt},
    major tick length=5pt,
    xlabel={$t$}, ylabel={$T$},
    xlabel style={anchor=west},
    ylabel style={anchor=south},
    xtick={-3,-2,-1,1,2,3},
    ytick={-5,-4,-3,-2,-1,1,2,3,4,5,6,7,8,9},
    tick label style={font=\tiny},
    label style={font=\scriptsize},
    samples=100,
]
\addplot[->, line width=1.25pt, domain=-2:4.5] {x^2-2*x-2};
\addplot[only marks, mark=*, mark size=2pt] coordinates {(-2,6) (1,-3)};
\end{axis}
\end{tikzpicture}
\end{center}

\begin{multipleChoice}
    \choice[correct]{Yes, $T$ is a function of $t$.}
    \choice{No, $T$ is a not a function of $t$.}
    \begin{feedback}
    \begin{solution}
domain =  $[-2, \infty)$ \\ range = $[-3, \infty)$
    \end{solution}
    \end{feedback}
\end{multipleChoice}
\end{problem}


\begin{problem}\label{graphfunctionfirsttT2}

    \begin{center}
\begin{tikzpicture}
\begin{axis}[
    width=2.7in, height=2.4in,
    xmin=-6, xmax=6, ymin=-6, ymax=6,
    axis lines=middle,
    axis line style={-{Latex[length=8pt,width=6pt]}},
    tick style={line width=1pt},
    major tick length=5pt,
    xlabel={$t$}, ylabel={$T$},
    xlabel style={anchor=west},
    ylabel style={anchor=south},
    xtick={-5,-4,-3,-2,-1,1,2,3,4,5},
    ytick={-5,-4,-3,-2,-1,1,2,3,4,5},
    tick label style={font=\tiny},
    label style={font=\scriptsize},
    samples=100,
]
% Polar curve r = 5 sin(3 theta)
\draw[line width=1.25pt,samples=400,domain=0:180,smooth,variable=\t]
  plot ({5*sin(3*\t)*cos(\t)},
        {5*sin(3*\t)*sin(\t)});
\end{axis}
\end{tikzpicture}
\end{center}




\begin{multipleChoice}
    \choice{Yes, $T$ is a function of $t$.}
    \choice[correct]{No, $T$ is a not a function of $t$.}
    \begin{feedback}
    The vertical line test fails in many places.
    \end{feedback}
\end{multipleChoice}
\end{problem}   



\begin{problem}\label{graphfunctionfirsttT3}
    Determine whether or not the graph suggests $T$ is a function of $t$  If it does, state the domain and range.

    \begin{center}
\begin{tikzpicture}
\begin{axis}[
    width=2.7in, height=2.4in,
    xmin=-6, xmax=6, ymin=-6, ymax=6,
    axis lines=middle,
    axis line style={-{Latex[length=8pt,width=6pt]}},
    tick style={line width=1pt},
    major tick length=5pt,
    xlabel={$t$}, ylabel={$T$},
    xlabel style={anchor=west},
    ylabel style={anchor=south},
    xtick={-5,-4,-3,-2,-1,1,2,3,4,5},
    ytick={-5,-4,-3,-2,-1,1,2,3,4,5},
    tick label style={font=\tiny},
    label style={font=\scriptsize},
    samples=100,
]
\addplot[line width=1.25pt, domain=-5:4] {0.0502*x^3-0.0344*x^2-0.2010*x+2.138};
\addplot[only marks, mark=*, mark size=2pt, mark options={fill=white}] coordinates {(-5,-4) (4,4)};
\end{axis}
\end{tikzpicture}
\end{center}

\begin{multipleChoice}
    \choice[correct]{Yes, $T$ is a function of $t$.}
    \choice{No, $T$ is a not a function of $t$.}
    \begin{feedback}
    \begin{solution}
domain =  $(-5, 4)$ \\ range = $(-4, 4)$
    \end{solution}
    \end{feedback}
\end{multipleChoice}
\end{problem}




\begin{problem}\label{graphfunctionlasttT}
    Determine whether or not the graph suggests $T$ is a function of $t$  If it does, state the domain and range.

    \begin{center}
\begin{tikzpicture}
\begin{axis}[
    width=2.7in, height=2.4in,
    xmin=-2, xmax=7, ymin=-6, ymax=6,
    axis lines=middle,
    axis line style={-{Latex[length=8pt,width=6pt]}},
    tick style={line width=1pt},
    major tick length=5pt,
    xlabel={$t$}, ylabel={$T$},
    xlabel style={anchor=west},
    ylabel style={anchor=south},
    xtick={-1,1,2,3,4,5,6},
    ytick={-5,-4,-3,-2,-1,1,2,3,4,5},
    tick label style={font=\tiny},
    label style={font=\scriptsize},
    samples=100,
]
\addplot[line width=1.25pt] coordinates {(0,-1) (3,-4)};
\addplot[line width=1.25pt] coordinates {(3,1) (4,4) (6,0)};
\addplot[only marks, mark=*, mark size=2pt] coordinates {(0,-1) (4,4) (6,0)};
\addplot[only marks, mark=*, mark size=2pt, mark options={fill=white}] coordinates {(3,-4) (3,1)};
\end{axis}
\end{tikzpicture}
\end{center}

\begin{multipleChoice}
    \choice[correct]{Yes, $T$ is a function of $t$.}
    \choice{No, $T$ is a not a function of $t$.}
    \begin{feedback}
    \begin{solution}
domain = $[0,3) \cup (3,6]$ \\ range = $(-4,-1] \cup [0,4]$
    \end{solution}
    \end{feedback}
\end{multipleChoice}
\end{problem} 
  
\end{question}

\begin{problem}
Determine which of the graphs, if any, in Exercises \ref{graphfunctionfirsttT} - \ref{graphfunctionlasttT}, suggest $t$ is a function of $T$.  (Feel free to use Exercise \ref{HLTExercise}.)  If it does, state the domain and range.

\begin{solution}
None of numbers \ref{graphfunctionfirsttT} - \ref{graphfunctionlasttT}  represent $t$ as a function of $T$.
\end{solution}
\end{problem}


%\setcounter{HW}{\value{enumi}}


\begin{question}
In Exercises \ref{graphfunctionfirstHs} - \ref{graphfunctionlastHs}, determine whether or not the graph suggests $H$ is a function of $s$.  For the ones which do, state the domain and range. 


\begin{problem}\label{graphfunctionfirstHs}

\begin{center}
\begin{tikzpicture}
\begin{axis}[
    width=2.7in, height=2.4in,
    xmin=-3, xmax=3, ymin=-1, ymax=5,
    axis lines=middle,
    axis line style={-{Latex[length=8pt,width=6pt]}},
    tick style={line width=1pt},
    major tick length=5pt,
    xlabel={$s$}, ylabel={$H$},
    xlabel style={anchor=west},
    ylabel style={anchor=south},
    xtick={-2,-1,1,2},
    ytick={1,2,3,4},
    tick label style={font=\tiny},
    label style={font=\scriptsize},
    samples=100,
]
\addplot[<->, line width=1.25pt, domain=-2.25:2.25] {4-x^2};
\end{axis}
\end{tikzpicture}
\end{center}

\begin{multipleChoice}
    \choice[correct]{Yes, $H$ is a function of $s$.}
    \choice{No, $H$ is a not a function of $s$.}
    \begin{feedback}
    \begin{solution}
domain = $(-\infty, \infty)$ \\ range = $(-\infty, 4]$
    \end{solution}
    \end{feedback}
\end{multipleChoice}
\end{problem}

\begin{problem}\label{graphfunctionfirstHs2}

\begin{center}
\begin{tikzpicture}
\begin{axis}[
    width=2.7in, height=2.4in,
    xmin=-3, xmax=3, ymin=-1, ymax=5,
    axis lines=middle,
    axis line style={-{Latex[length=8pt,width=6pt]}},
    tick style={line width=1pt},
    major tick length=5pt,
    xlabel={$s$}, ylabel={$H$},
    xlabel style={anchor=west},
    ylabel style={anchor=south},
    xtick={-2,-1,1,2},
    ytick={1,2,3,4},
    tick label style={font=\tiny},
    label style={font=\scriptsize},
    samples=100,
]
\addplot[<->, line width=1.25pt] coordinates {(-2,-1) (1,4) (2,-1)};
\end{axis}
\end{tikzpicture}
\end{center}

\begin{multipleChoice}
    \choice[correct]{Yes, $H$ is a function of $s$.}
    \choice{No, $H$ is a not a function of $s$.}
    \begin{feedback}
    \begin{solution}
domain = $(-\infty, \infty)$ \\ range = $(-\infty, 4]$
    \end{solution}
    \end{feedback}
\end{multipleChoice}
\end{problem}


\begin{problem}\label{graphfunctionfirstHs3}

\begin{center}
\begin{tikzpicture}
\begin{axis}[
    width=2.7in, height=2.4in,
    xmin=-3, xmax=3, ymin=-1, ymax=5,
    axis lines=middle,
    axis line style={-{Latex[length=8pt,width=6pt]}},
    tick style={line width=1pt},
    major tick length=5pt,
    xlabel={$s$}, ylabel={$H$},
    xlabel style={anchor=west},
    ylabel style={anchor=south},
    xtick={-2,-1,1,2},
    ytick={1,2,3,4},
    tick label style={font=\tiny},
    label style={font=\scriptsize},
    samples=100,
]
\addplot[->, line width=1.25pt, domain=-2:1.8] {3-2*sqrt(x+2)};
\addplot[only marks, mark=*, mark size=2pt] coordinates {(-2,3)};
\end{axis}
\end{tikzpicture}
\end{center}

\begin{multipleChoice}
    \choice[correct]{Yes, $H$ is a function of $s$.}
    \choice{No, $H$ is a not a function of $s$.}
    \begin{feedback}
    \begin{solution}
domain =  $[-2, \infty)$  \\ range =   $(-\infty, 3]$
    \end{solution}
    \end{feedback}
\end{multipleChoice}
\end{problem}

\begin{problem}\label{graphfunctionlastHs}


\begin{center}
\begin{tikzpicture}
\begin{axis}[
    width=2.7in, height=2.4in,
    xmin=-3, xmax=3, ymin=-1, ymax=5,
    axis lines=middle,
    axis line style={-{Latex[length=8pt,width=6pt]}},
    tick style={line width=1pt},
    major tick length=5pt,
    xlabel={$s$}, ylabel={$H$},
    xlabel style={anchor=west},
    ylabel style={anchor=south},
    xtick={-2,-1,1,2},
    ytick={1,2,3,4},
    tick label style={font=\tiny},
    label style={font=\scriptsize},
    samples=100,
]
\addplot[<->, line width=1.25pt, domain=-2.15:1.75] {x*(x-1)*(x+2)};
\end{axis}
\end{tikzpicture}
\end{center}

\begin{multipleChoice}
    \choice[correct]{Yes, $H$ is a function of $s$.}
    \choice{No, $H$ is a not a function of $s$.}
    \begin{feedback}
    \begin{solution}
domain = $(-\infty, \infty)$ \\ range = $(-\infty, \infty)$
    \end{solution}
    \end{feedback}
\end{multipleChoice}
\end{problem}

\end{question}

\begin{problem}
Determine which, if any, of the graphs in numbers \ref{graphfunctionfirstHs} - \ref{graphfunctionlastHs} represent $s$ as a function of $H$.  For the ones which do, state the domain and range.   (Feel free to use Exercise \ref{HLTExercise}.)

\begin{solution}
Only number \ref{graphfunctionfirstHs3} represents $s$ as a function of $H$;  domain =  $(-\infty, 3]$  and range =   $[-2, \infty)$
\end{solution}
\end{problem} 

\begin{question}
In Exercises \ref{graphfunctionfirstut} - \ref{graphfunctionlastut}, determine whether or not the graph suggests $u$ is a function of $t$. For the ones which do, state the domain and range. 

\begin{problem}\label{graphfunctionfirstut}

\begin{center}
\begin{tikzpicture}
\begin{axis}[
    width=2.7in, height=2.4in,
    xmin=-3, xmax=3, ymin=-3, ymax=3,
    axis lines=middle,
    axis line style={-{Latex[length=8pt,width=6pt]}},
    tick style={line width=1pt},
    major tick length=5pt,
    xlabel={$t$}, ylabel={$u$},
    xlabel style={anchor=west},
    ylabel style={anchor=south},
    xtick={-2,-1,1,2},
    ytick={-2,-1,1,2},
    tick label style={font=\tiny},
    label style={font=\scriptsize},
    samples=100,
]
\addplot[->, line width=1.25pt] coordinates {(0,1) (-2,-2)};
\addplot[->, line width=1.25pt] coordinates {(1,2) (3,2)};
\addplot[only marks, mark=*, mark size=2pt] coordinates {(0,1)};
\addplot[only marks, mark=*, mark size=2pt, mark options={fill=white}] coordinates {(1,2)};
\end{axis}
\end{tikzpicture}
\end{center}

\begin{multipleChoice}
    \choice[correct]{Yes, $u$ is a function of $t$.}
    \choice{No, $u$ is a not a function of $t$.}
    \begin{feedback}
    \begin{solution}
domain =  $(-\infty, 0] \cup (1, \infty)$ \\ range =  $(-\infty, 1] \cup \{ 2\}$
    \end{solution}
    \end{feedback}
\end{multipleChoice}
\end{problem} 

\begin{problem}\label{graphfunctionfirstut2}


\begin{center}
\begin{tikzpicture}
\begin{axis}[
    width=2.7in, height=2.4in,
    xmin=-4, xmax=4, ymin=-3, ymax=3,
    axis lines=middle,
    axis line style={-{Latex[length=8pt,width=6pt]}},
    tick style={line width=1pt},
    major tick length=5pt,
    xlabel={$t$}, ylabel={$u$},
    xlabel style={anchor=west},
    ylabel style={anchor=south},
    xtick={-3,-2,-1,1,2,3},
    ytick={-2,-1,1,2},
    tick label style={font=\tiny},
    label style={font=\scriptsize},
    samples=100,
]
\addplot[line width=1.25pt, domain=-3:3] {2*sin(deg(1.05*x))};
\addplot[only marks, mark=*, mark size=2pt] coordinates {(-3,0) (3,0)};
\end{axis}
\end{tikzpicture}
\end{center}

\begin{multipleChoice}
    \choice[correct]{Yes, $u$ is a function of $t$.}
    \choice{No, $u$ is a not a function of $t$.}
    \begin{feedback}
    \begin{solution}
domain =  $[-3,3]$ \\ range =  $[-2,2]$
    \end{solution}
    \end{feedback}
\end{multipleChoice}
\end{problem} 

\begin{problem}\label{graphfunctionfirstut3}

\begin{center}
\begin{tikzpicture}
\begin{axis}[
    width=2.7in, height=2.4in,
    xmin=-3, xmax=3, ymin=-3, ymax=3,
    axis lines=middle,
    axis line style={-{Latex[length=8pt,width=6pt]}},
    tick style={line width=1pt},
    major tick length=5pt,
    xlabel={$t$}, ylabel={$u$},
    xlabel style={anchor=west},
    ylabel style={anchor=south},
    xtick={-2,-1,1,2},
    ytick={-2,-1,1,2},
    tick label style={font=\tiny},
    label style={font=\scriptsize},
    samples=100,
]
\addplot[<->, line width=1.25pt] coordinates {(2,-3) (2,3)};
\end{axis}
\end{tikzpicture}
\end{center}

\begin{multipleChoice}
    \choice{Yes, $u$ is a function of $t$.}
    \choice[correct]{No, $u$ is a not a function of $t$.}
    \begin{feedback}
    The vertical line test fails at $t=2$.
    \end{feedback}
\end{multipleChoice}
\end{problem}  

\begin{problem}\label{graphfunctionlastut}


\begin{center}
\begin{tikzpicture}
\begin{axis}[
    width=2.7in, height=2.4in,
    xmin=-3, xmax=3, ymin=-3, ymax=3,
    axis lines=middle,
    axis line style={-{Latex[length=8pt,width=6pt]}},
    tick style={line width=1pt},
    major tick length=5pt,
    xlabel={$t$}, ylabel={$u$},
    xlabel style={anchor=west},
    ylabel style={anchor=south},
    xtick={-2,-1,1,2},
    ytick={-2,-1,1,2},
    tick label style={font=\tiny},
    label style={font=\scriptsize},
    samples=100,
]
\addplot[<->, line width=1.25pt] coordinates {(-3,2) (3,2)};
\end{axis}
\end{tikzpicture}
\end{center}

\begin{multipleChoice}
    \choice[correct]{Yes, $u$ is a function of $t$.}
    \choice{No, $u$ is a not a function of $t$.}
    \begin{feedback}
    \begin{solution}
domain = $(-\infty, \infty)$ \\ range = $\{2\}$
    \end{solution}
    \end{feedback}
\end{multipleChoice}
\end{problem} 
\end{question}

\begin{problem}
Determine which, if any, of the graphs in numbers \ref{graphfunctionfirstut} - \ref{graphfunctionlastut} represent $t$ as a function of $u$.  For the ones which do, state the domain and range.   (Feel free to use Exercise \ref{HLTExercise}.)

\begin{solution}
Only number \ref{graphfunctionfirstut3} represents $t$ as a function of $u$;  domain = $(-\infty, \infty)$ and range=$\{2 \}$.
\end{solution}
\end{problem} 

\begin{question}
In Exercises \ref{functionvaluesfromgraphfirst} - \ref{functionvaluesfromgraphlast}, use the graphs of $f$ and $g$ below to find the indicated values.



\begin{tikzpicture}
\begin{axis}[
    width=2.7in, height=3in,
    xmin=-6, xmax=6, ymin=-6, ymax=6,
    axis lines=middle,
    axis line style={-{Latex[length=8pt,width=6pt]}},
    tick style={line width=1pt},
    major tick length=5pt,
    xlabel={$x$}, ylabel={$y$},
    xlabel style={anchor=west},
    ylabel style={anchor=south},
    xtick={-5,-4,-3,-2,-1,1,2,3,4,5},
    ytick={-5,-4,-3,-2,-1,1,2,3,4,5},
    tick label style={font=\tiny},
    label style={font=\scriptsize},
    samples=100,
]
\addplot[only marks, mark=*, mark size=2pt] coordinates {(-5,-5) (-4,0) (-3,4) (-2,2) (-1,0) (0,-1) (1,0) (2,3) (3,1)};
\addplot[line width=1.25pt] coordinates {(-5,-5) (-4,0) (-3,4) (-2,2) (-1,0)};
\addplot[line width=1.25pt, domain=-1:1] {x^2-1};
\addplot[line width=1.25pt] coordinates {(1,0) (2,3) (3,1)};
\end{axis}
\node at (rel axis cs:0.5,0) [below] {The graph of $y=f(x)$.};
\end{tikzpicture}
\begin{tikzpicture}
\begin{axis}[
    width=2.7in, height=3in,
    xmin=-5, xmax=5, ymin=-6, ymax=6,
    axis lines=middle,
    axis line style={-{Latex[length=8pt,width=6pt]}},
    tick style={line width=1pt},
    major tick length=5pt,
    xlabel={$t$}, ylabel={$s$},
    xlabel style={anchor=west},
    ylabel style={anchor=south},
    xtick={-4,-3,-2,-1,1,2,3,4},
    ytick={-5,-4,-3,-2,-1,1,2,3,4,5},
    tick label style={font=\tiny},
    label style={font=\scriptsize},
    samples=100,
]
\addplot[line width=1.25pt, domain=-4:4] {5*sin(deg(pi*x/4))};
\addplot[only marks, mark=*, mark size=2pt] coordinates {(-2,-5) (0,0) (4,0) (2,3) (-4,0)};
\addplot[only marks, mark=*, mark size=2pt, mark options={fill=white}] coordinates {(2,5)};
\end{axis}
\node at (rel axis cs:0.5,0) [below] {The graph of $s=g(t)$.};
\end{tikzpicture}




\begin{problem}\label{functionvaluesfromgraphfirst} 
$f(-2)=\answer{2}$
\end{problem}

\begin{problem}
$g(-2)=\answer{-5}$
\end{problem}

\begin{problem}
$f(2)=\answer{3}$
\end{problem}

\begin{problem}
$g(2)=\answer{3}$
\end{problem}

\begin{problem}
$f(0)=\answer{-1}$
\end{problem}

\begin{problem}
$g(0)=\answer{0}$
\end{problem}

\begin{problem}
Solving $f(x)=0$ gives us 
\begin{selectAll}
    \choice[correct]{$x=-4$}
    \choice{$x=-3$}
    \choice[correct]{$x=-1$}
    \choice{$x=0$}
    \choice{$x=4$}
\end{selectAll}
\end{problem}

\begin{problem}
Solving $g(t)=0$ gives us 
\begin{selectAll}
    \choice[correct]{$t=-4$}
    \choice{$t=-2$}
    \choice[correct]{$t=0$}
    \choice{$t=2$}
    \choice[correct]{$t=4$}
\end{selectAll}
\end{problem}

\begin{problem}
State the domain and range of $f$.

\begin{solution}
Domain:  $[-5,3]$,  Range:  $[-5,4]$.
\end{solution}
\end{problem}

\begin{problem}\label{functionvaluesfromgraphlast}
State the domain and range of $g$ .

\begin{solution}
Domain:  $[-4,4]$,  Range:  $[-5,5)$.
\end{solution}
\end{problem}
\end{question}

\begin{question}
In Exercises \ref{sketchgraphfirst} - \ref{sketchgraphlast}, graph each function by making a table, plotting points, and using a graphing utility (if needed.)  Use the independent variable as the horizontal axis label and the default `$y$' label for the vertical axis label.  State the domain and range of each function. 

\begin{problem}\label{sketchgraphfirst}
$f(x) = 2-x$
\end{problem}

\begin{problem}
$g(t) = \dfrac{t - 2}{3}$
\end{problem} 

\begin{problem}
$f(x) = 4-x^2$
\end{problem} 

\begin{problem}
$g(t) = 2$
\end{problem} 

\begin{problem}
$h(s) = s^3$
\end{problem} 

\begin{problem}
$f(x) = x(x-1)(x+2)$
\end{problem} 

\begin{problem}
$g(t) = \sqrt{t-2}$
\end{problem} 

\begin{problem}
$h(s) = \sqrt{5 - s}$
\end{problem}

\begin{problem}
$f(x) = 3-2\sqrt{x+2}$
\end{problem}

\begin{problem}
$g(t) = \sqrt[3]{t}$
\end{problem}

\begin{problem}\label{sketchgraphlast}
$h(s) = \dfrac{1}{s^{2} + 1}$
\end{problem}

\end{question}



\begin{problem}
Consider the function $f$ described below:

\begin{center}
\begin{tikzpicture}[>=Latex, scale=.75]
\node at (-4.25,5) {$-1$};
\node at (-4.25,3) {0};
\node at (-4.25,1) {1};
\node at (-4.25,-1) {2};
\node at (0,6) {$f$};
\node at (3.5,5) {$-3$};
\node at (3.5,3) {0};
\node at (3.5,1) {4};
\draw[->, line width=1pt] (-2.5,5) -- (2.5,3);
\draw[->, line width=1pt] (-2.5,3) -- (2.5,5);
\draw[->, line width=1pt] (-2.5,1) -- (2.5,3);
\draw[->, line width=1pt] (-2.5,-1) -- (2.5,1);
\end{tikzpicture}
\end{center}
% \end{center}

\begin{enumerate}

\item  State  the domain and range of $f$.

\item Find $f(0)$ and solve $f(x) = 0$.

\item  Write $f$ as a set of ordered pairs.

\item  Graph $f$.

\end{enumerate}

\end{problem}

\begin{problem}
Let $g = \{ (-1,4), (0,2), (2, 3), (3,4)  \}$

\begin{enumerate}

\item  State the domain and range of $g$.

\item  Create a mapping diagram for $g$.

\item  Find $g(0)$ and solve $g(x) = 0$.

\item  Graph $g$.

\end{enumerate}
\end{problem}

\begin{problem}
Let $F = \{ (t, t^2) \, | \, \text{$t$ is a real number} \}$.  Find $F(4)$ and solve $F(x) = 4$.

\begin{hint}
Elements of $F$ are of the form $(x, F(x))$.
\end{hint}
\end{problem}

\begin{problem}
Let $G = \{ (2t, t+5) \, | \, \text{$t$ is a real number} \}$.  Find $G(4)$ and solve $G(x) = 4$.

\begin{hint}
Elements of $G$ are of the form $(x, G(x))$.
\end{hint}
\end{problem}

\begin{problem}
The area enclosed by a square, in square inches,  is a function of the length of one of its sides $\ell$, when measured in inches.  This function is represented by the formula $A(\ell) = \ell^2$ for $\ell > 0$.  Find $A(3)$ and solve $A(\ell) = 36$.  Interpret your answers to each.  Why is $\ell$ restricted to $\ell > 0$?
\end{problem}

\begin{problem}
The area enclosed by a circle, in square meters, is a function of its radius $r$, when measured in meters.  This function is represented by the formula $A(r) = \pi r^2$ for $r > 0$.  Find $A(2)$ and solve $A(r) = 16\pi$.  Interpret your answers to each.  Why is $r$ restricted to $r > 0$?
\end{problem}

\begin{problem}
The volume  enclosed by a cube, in cubic centimeters, is a function of the length of one of its sides $s$, when measured in centimeters.   This function is represented by the formula $V(s) = s^3$ for $s > 0$.  Find $V(5)$ and solve $V(s) = 27$.  Interpret your answers to each.  Why is $s$ restricted to $s > 0$?
\end{problem}

\begin{problem}
The volume enclosed by a sphere, in cubic feet, is a function of the radius of the sphere $r$, when measured in feet. This function is represented by the formula  $V(r) =\frac{4\pi}{3} r^{3}$ for $r > 0$.  Find $V(3)$ and solve $V(r) = \frac{32\pi}{3}$.  Interpret your answers to each.  Why is $r$ restricted to $r > 0$?
\end{problem}

\begin{problem}
The height of an object dropped from the roof of an eight story building is modeled by the function:  $h(t) = -16t^2 + 64$, $0 \leq t \leq 2$. Here,  $h(t)$ is the height of the object off the ground, in feet, $t$ seconds after the object is dropped.  Find $h(0)$ and solve $h(t) = 0$.  Interpret your answers to each.  Why is $t$ restricted to $0 \leq t \leq 2$?
\end{problem}

\begin{problem}
The temperature in degrees Fahrenheit $t$ hours after 6 AM is given by $T(t) = -\frac{1}{2} t^2 + 8t+3$ for $0 \leq t \leq 12$. Find and interpret $T(0)$, $T(6)$ and $T(12)$. 
\end{problem}

\begin{problem}
The function $C(x) = x^2-10x+27$  models the cost, in \textit{hundreds} of dollars, to produce $x$ \textit{thousand} pens.  Find and interpret $C(0)$, $C(2)$ and $C(5)$.
\end{problem}

\begin{problem}
Using data from the  \href{http://www.bts.gov/publications/national_transportation_statistics/html/table_04_23.html}{\underline{Bureau of Transportation Statistics}}, the average fuel economy in miles per gallon for passenger cars in the US can be modeled by  $E(t) = -0.0076t^2+0.45t + 16$, $0 \leq t \leq 28$, where $t$ is the number of years since $1980$. Use a calculator to find $E(0)$, $E(14)$ and $E(28)$.  Round your answers to two decimal places and interpret your answers to each.
\end{problem} 

\begin{problem}
The perimeter of a square, in centimeters,  is four times the length of one if its sides, also measured in centimeters.  Represent the function $P$ which takes as its input the length of the side of a square in centimeters, $s$ and returns the perimeter of the square in inches, $P(s)$ using a formula.
\end{problem}

\begin{problem}
The circumference of a circle, in feet,  is $\pi$ times the diameter of the circle, also measured in feet.  Represent the function $C$ which takes as its input the length of the diameter of a circle in feet, $D$ and returns the circumference of a circle in inches, $C(D)$ using a formula.
\end{problem}

\begin{problem}
Suppose $A(P)$ gives the amount of money in a retirement account (in dollars) after 30 years as a function of the amount of the monthly payment (in dollars), $P$.

\begin{enumerate}

\item What does $A(50)$ mean?  

\item  What is the significance of the solution to the equation $A(P) = 250000$? .

\item  Explain what each of the following expressions mean:  $A(P+50)$, $A(P)+50$, and $A(P) + A(50)$.  

\end{enumerate}
\end{problem}
 
\begin{problem}
Suppose $P(t)$ gives the chance of precipitation (in percent)  $t$ hours after 8 AM.

\begin{enumerate}

\item  Write an expression which gives the chance of precipitation at noon.

\item  Write an inequality which determines when the chance of precipitation is more than $50 \%$.

\end{enumerate}
\end{problem}

\begin{problem}
Explain why the graph in  Exercise \ref{graphfunctionfirstvw}  suggests that not only is $v$ as a function of $w$ but also $w$ is a function of $v$.  Suppose $w = f(v)$ and $v = g(w)$.  That is, $f$ is the name of the function which takes $v$ values as inputs and returns $w$ values as outputs and $g$ is the name of the function which does vice-versa.   Find the domain and range of  $g$ and compare these to the domain and range of $f$.
\end{problem}

\begin{problem}
Sketch the graph of a function with domain $(-\infty, 3) \cup [4,5)$ with range $\{ 2 \} \cup (5, \infty)$.
\end{problem}

  


\end{document}
